\documentclass[../DoAn.tex]{subfiles}
\begin{document}

Chương này trình bày thiết kế và quá trình xây dựng ứng dụng. Nội dung chương mô tả kiến trúc phần mềm, thiết kế các thành phần chính, giao diện, dữ liệu lưu trữ, công cụ sử dụng, kết quả đạt được và kế hoạch kiểm thử.

\section{Thiết kế kiến trúc}
\subsection{Lựa chọn kiến trúc phần mềm}
Ứng dụng được thiết kế theo hướng phân tách trách nhiệm giữa dữ liệu, logic nghiệp vụ và giao diện. Do Godot tổ chức chương trình theo scene và node, kiến trúc của đồ án không áp dụng MVC theo hình thức truyền thống, nhưng vẫn giữ tinh thần tách lớp: thành phần trạng thái trò chơi lưu dữ liệu và quy tắc tính toán, thành phần sinh màn chịu trách nhiệm tạo dữ liệu chơi, thành phần giao diện chịu trách nhiệm hiển thị và nhận tương tác từ người dùng.

Trong kiến trúc cụ thể, \texttt{GameState} quản lý bảng chữ, danh sách từ đã đặt, điểm số, thời gian, lịch sử đoán và nhiệm vụ. \texttt{BoardGenerator} sinh bảng ô chữ từ danh sách từ. \texttt{ThemedLevelProvider} quản lý chủ đề, gợi ý và chia chủ đề thành các màn. \texttt{UIController} điều khiển màn hình chơi, còn \texttt{main.gd} điều khiển luồng menu, tài khoản và chuyển cảnh. \texttt{SaveManager} quản lý việc đọc ghi tiến độ cục bộ. Cách phân tách này giúp từng phần có nhiệm vụ tương đối rõ ràng và có thể thay đổi độc lập ở mức nhất định.

Luồng phụ thuộc được thiết kế theo hướng các thành phần giao diện gọi tới các thành phần nghiệp vụ, trong khi các thành phần nghiệp vụ không phụ thuộc ngược vào giao diện cụ thể. Ví dụ, \texttt{BoardGenerator} không biết bảng được hiển thị bằng nút hay bằng hình ảnh; nó chỉ trả về cấu trúc lưới và danh sách từ đã đặt. \texttt{GameState} không biết dữ liệu được lấy từ chủ đề nào; nó chỉ nhận một bảng đã sinh và quản lý trạng thái chơi trên bảng đó. Cách tách này giúp việc kiểm tra logic sinh bảng và logic trạng thái không bị ràng buộc chặt với scene.

\begin{table}[H]
\centering
\begin{tabularx}{\textwidth}{|>{\RaggedRight\arraybackslash}p{3.2cm}|>{\RaggedRight\arraybackslash}p{5.2cm}|>{\RaggedRight\arraybackslash}X|}
\hline
\textbf{Tầng/nhóm} & \textbf{Thành phần} & \textbf{Trách nhiệm chính} \\ \hline
Điều phối ứng dụng & \classname{main.gd} & Quản lý luồng menu, tài khoản, chọn chủ đề, chọn màn và chuyển sang màn chơi. \\ \hline
Giao diện chơi & \makecell[l]{\classname{UIController}\\\classname{TileButton}\\\classname{InputController}} & Hiển thị bảng, gợi ý, trạng thái, nhận sự kiện từ người chơi và cập nhật giao diện. \\ \hline
Trạng thái trò chơi & \classname{GameState} & Lưu bảng, từ đã đặt, điểm, thời gian, nhiệm vụ, lịch sử đoán và trạng thái hoàn thành. \\ \hline
Sinh nội dung & \makecell[l]{\classname{ThemedLevelProvider}\\\classname{BoardGenerator}} & Quản lý chủ đề, chia màn, tạo gợi ý và sinh bảng ô chữ. \\ \hline
Xử lý tiếng Việt & \makecell[l]{\classname{WordEntry}\\\classname{VietnameseNormalizer}\\\classname{DiacriticHelper}} & Chuẩn hóa từ, tách ký tự, quản lý dấu và so khớp đáp án. \\ \hline
Lưu trữ & \classname{SaveManager} & Đọc, ghi, chuẩn hóa dữ liệu tiến độ cục bộ theo tài khoản. \\ \hline
\end{tabularx}
\caption{Phân nhóm thành phần trong kiến trúc ứng dụng}
\label{table:architecture_groups}
\end{table}

\section{Thiết kế tổng quan}
\thesistodo{vẽ biểu đồ kiến trúc/gói UML cho hệ thống. Biểu đồ nên có các gói hoặc nhóm thành phần: Menu, Gameplay UI, Game State, Level Provider, Board Generator, Save Manager, Vietnamese Processing. Mũi tên phụ thuộc nên thể hiện Main gọi UIController và ThemedLevelProvider; UIController dùng GameState, InputController và SaveManager; ThemedLevelProvider dùng WordEntry, VietnameseNormalizer và BoardGenerator.}

Ở mức tổng quan, hệ thống bắt đầu từ scene \texttt{Main}. Scene này chứa \texttt{StartMenu} và \texttt{GameBoard}. Khi người chơi chọn một màn, \texttt{main.gd} yêu cầu \texttt{ThemedLevelProvider} tạo dữ liệu màn, sau đó khởi tạo \texttt{GameState} và truyền cho \texttt{UIController}. Từ thời điểm đó, \texttt{UIController} chịu trách nhiệm xây dựng bảng, hiển thị gợi ý, nhận đáp án, cập nhật trạng thái và gọi \texttt{SaveManager} khi cần lưu dữ liệu.

Scene \texttt{StartMenu} chỉ chứa khung giao diện tĩnh ban đầu như nền, thẻ menu, hàng chọn tài khoản, hàng chọn chủ đề và các nút điều hướng. Các popup như quản lý tài khoản, thêm tài khoản, xác nhận xóa tài khoản, hướng dẫn và chọn màn được tạo bằng mã trong \texttt{main.gd}. Cách này giúp giảm số lượng scene phụ, đồng thời cho phép nội dung popup được cập nhật theo trạng thái hiện tại như danh sách tài khoản hoặc danh sách màn.

Scene \texttt{GameBoard} chứa cấu trúc giao diện của màn chơi. Bảng chữ được đặt trong \texttt{GridContainer}, danh sách gợi ý được đặt trong \texttt{HFlowContainer}, còn các nút thao tác như chơi lại, mở gợi ý và đoán đáp án được đặt trong vùng điều khiển. Các overlay tạm dừng và hoàn thành cũng nằm trong scene này nhưng mặc định ẩn, chỉ hiện khi người chơi bấm phím tạm dừng hoặc hoàn thành toàn bộ màn chơi.

\section{Thiết kế chi tiết gói/module}
\thesistodo{vẽ biểu đồ lớp mức cao. Nên thể hiện các lớp chính gồm BoardGenerator, WordEntry, VietnameseNormalizer, DiacriticHelper, ThemedLevelProvider, GameState, UIController, InputController, SaveManager và TileButton.}

Nhóm xử lý dữ liệu từ gồm \texttt{WordEntry}, \texttt{VietnameseNormalizer} và \texttt{DiacriticHelper}. \texttt{WordEntry} giữ dạng từ gốc, dạng xếp bảng, gợi ý, danh sách ký tự và độ dài. \texttt{VietnameseNormalizer} cung cấp các hàm chuẩn hóa, bỏ khoảng trắng và so khớp đáp án. \texttt{DiacriticHelper} mô tả thông tin dấu và biến thể nguyên âm của từng ký tự.

Nhóm sinh nội dung gồm \texttt{ThemedLevelProvider} và \texttt{BoardGenerator}. \texttt{ThemedLevelProvider} tạo danh sách từ theo chủ đề, sinh gợi ý và chia chủ đề thành nhiều màn. \texttt{BoardGenerator} nhận danh sách \texttt{WordEntry} và sinh ra lưới ô chữ cùng danh sách các từ đã đặt. Nhóm giao diện gồm \texttt{main.gd}, \texttt{UIController}, \texttt{InputController} và \texttt{TileButton}. Các thành phần này nhận sự kiện từ người chơi và cập nhật hiển thị dựa trên trạng thái trò chơi.

\begin{table}[H]
\centering
\begin{tabularx}{\textwidth}{|>{\RaggedRight\arraybackslash}p{4.6cm}|>{\RaggedRight\arraybackslash}X|>{\RaggedRight\arraybackslash}X|}
\hline
\textbf{Lớp/script} & \textbf{Dữ liệu chính} & \textbf{Phương thức/chức năng tiêu biểu} \\ \hline
\classname{BoardGenerator} & Lưới tạm, danh sách từ đã đặt & Sinh bảng, tìm vị trí đặt từ, kiểm tra hợp lệ, tính giao điểm \\ \hline
\classname{ThemedLevelProvider} & Danh sách chủ đề, gợi ý, màn chơi & Xây dựng chủ đề, chia màn, tạo dữ liệu màn, sinh gợi ý \\ \hline
\classname{GameState} & Bảng chữ, từ đã đặt, điểm, thời gian, nhiệm vụ & Kiểm tra từ, mở từ, tính điểm, quản lý nhiệm vụ, đếm tiến độ \\ \hline
\classname{UIController} & Node giao diện, ô chữ, nút gợi ý, trạng thái hiển thị & Dựng bảng, nhận đáp án, xử lý gợi ý, cập nhật UI, hiển thị kết quả \\ \hline
\classname{SaveManager} & Hồ sơ người chơi, bản ghi tiến độ & Tạo tài khoản, chọn tài khoản, lưu điểm, lưu sao, lưu tiến độ màn \\ \hline
\classname{VietnameseNormalizer} & Chuỗi tiếng Việt & Chuẩn hóa từ, bỏ khoảng trắng, tách ký tự, so khớp đáp án \\ \hline
\end{tabularx}
\caption{Mô tả các lớp và script chính}
\label{table:main_classes}
\end{table}

\section{Thiết kế giao diện}
Giao diện được chia thành hai phần chính: menu và màn chơi. Menu có nhiệm vụ chọn tài khoản, chọn chủ đề và chọn màn. Màn chơi hiển thị bảng ô chữ ở bên trái, các gợi ý và nút thao tác ở bên phải, phía trên có thông tin thời gian và số sao. Trên màn hình hẹp, bố cục chuyển sang dạng dọc để hạn chế tràn nội dung.

\thesistodo{chèn ảnh màn hình menu chính, màn chọn chủ đề/màn chơi và màn chơi. Ảnh nên lưu trong thư mục \texttt{Hinhve/}, sau đó dùng \texttt{\textbackslash includegraphics}.}

Màn hình menu chính được thiết kế để người chơi bắt đầu nhanh. Thành phần nổi bật nhất là tiêu đề trò chơi, hàng chọn tài khoản và các nút Chơi, Hướng dẫn, Thoát game. Khi bấm Chơi, giao diện chuyển sang chế độ chọn chủ đề. Ở chế độ này, hàng tài khoản được ẩn, hàng chủ đề được hiển thị và nút quay lại xuất hiện ở góc trên. Cách tách luồng này giúp màn hình không bị quá nhiều lựa chọn cùng lúc.

Popup chọn màn xuất hiện sau khi người chơi chọn một chủ đề và bấm chọn màn. Mỗi nút màn hiển thị số thứ tự màn, số từ trong màn, số sao đã đạt và trạng thái đã hoàn thành hoặc chưa hoàn thành. Thiết kế này giúp người chơi biết nhanh tiến độ của từng chủ đề mà không cần mở từng màn.

Trong màn chơi, vùng bảng chữ chiếm phần lớn chiều ngang vì đây là khu vực tương tác chính. Vùng điều khiển bên phải chứa danh sách số gợi ý. Khi người chơi chọn một số, nội dung gợi ý được hiển thị ở tiêu đề vùng gợi ý. Cách hiển thị này giúp danh sách gợi ý không chiếm quá nhiều diện tích, đặc biệt khi màn có nhiều từ. Ô nhập đáp án và nút đoán được đặt gần vùng gợi ý để người chơi có thể đọc gợi ý và nhập câu trả lời liên tục.

Giao diện cũng có các phản hồi trạng thái như thông báo đúng, sai, không đủ sao hoặc yêu cầu chọn gợi ý. Các phản hồi này được hiển thị trong vùng trạng thái thay vì dùng nhiều hộp thoại nhỏ, nhằm giảm gián đoạn khi chơi. Khi người chơi hoàn thành toàn bộ màn, hộp thoại kết quả mới được hiển thị để tổng kết điểm, sao, thời gian và nhiệm vụ.

\section{Thiết kế lớp}
\texttt{GameState} là lớp trung tâm của trạng thái chơi. Lớp này lưu lưới ô chữ, danh sách từ, số lần đoán, số đáp án đúng, điểm, nhiệm vụ, thời gian, gợi ý đã dùng và lịch sử nhập đáp án. Các phương thức quan trọng gồm đặt bảng, kiểm tra từ, mở từ, mở một chữ, tính điểm, tính nhiệm vụ và lấy thời gian đã chơi.

\texttt{BoardGenerator} chịu trách nhiệm sinh bảng. Phương thức chính \texttt{generate\_board} tạo lưới rỗng, đặt từ đầu tiên, sau đó tìm vị trí phù hợp cho các từ còn lại. Các hàm phụ như tìm vị trí giao nhau, kiểm tra khả năng đặt từ, tính số giao điểm và tính vùng bao giúp đánh giá chất lượng vị trí đặt từ.

\texttt{UIController} là lớp điều khiển giao diện màn chơi. Lớp này xây dựng các ô chữ từ \texttt{Tile.tscn}, tạo danh sách nút gợi ý, nhận sự kiện nhập đáp án, xử lý nút gợi ý, cập nhật điểm, sao, thời gian và hiển thị hộp thoại hoàn thành. \texttt{SaveManager} là lớp tĩnh dùng để lưu và đọc dữ liệu tiến độ từ tệp JSON.

\begin{table}[H]
\centering
\begin{tabularx}{\textwidth}{|>{\RaggedRight\arraybackslash}p{4.6cm}|>{\RaggedRight\arraybackslash}X|}
\hline
\textbf{Lớp} & \textbf{Thiết kế chi tiết} \\ \hline
\classname{GameState} & Quản lý toàn bộ dữ liệu động của một lượt chơi. Lớp này không trực tiếp hiển thị giao diện mà cung cấp các hàm như kiểm tra từ, đếm số từ hoàn thành, mở từ, mở một chữ và tính nhiệm vụ. \\ \hline
\classname{BoardGenerator} & Là lớp thuật toán. Lớp này nhận danh sách \classname{WordEntry}, chiều rộng và chiều cao bảng, sau đó trả về từ điển gồm lưới và danh sách từ đã đặt. \\ \hline
\classname{ThemedLevelProvider} & Là lớp dữ liệu và chuẩn bị màn chơi. Lớp này chứa danh sách chủ đề, tạo \classname{WordEntry}, chia màn, sinh gợi ý và gọi bộ sinh bảng khi người chơi bắt đầu một màn. \\ \hline
\classname{UIController} & Là lớp trung gian giữa người chơi và \classname{GameState}. Lớp này biến dữ liệu trạng thái thành giao diện, đồng thời chuyển thao tác của người chơi thành lời gọi tới trạng thái trò chơi. \\ \hline
\classname{SaveManager} & Được thiết kế dưới dạng các hàm tĩnh để các lớp khác có thể đọc ghi tiến độ mà không cần giữ một instance riêng. \\ \hline
\end{tabularx}
\caption{Thiết kế chi tiết các lớp chủ đạo}
\label{table:class_design}
\end{table}

\thesistodo{vẽ 1 hoặc 2 biểu đồ trình tự. Gợi ý biểu đồ 1: người chơi chọn màn -> Main -> ThemedLevelProvider -> BoardGenerator -> GameState -> UIController. Gợi ý biểu đồ 2: người chơi nhập đáp án -> UIController -> VietnameseNormalizer -> GameState -> SaveManager nếu hoàn thành.}

\section{Thiết kế dữ liệu lưu trữ}
Ứng dụng không sử dụng cơ sở dữ liệu quan hệ. Dữ liệu tiến độ được lưu dưới dạng JSON trong tệp \texttt{user://records.json}. Cấu trúc dữ liệu gồm hồ sơ người chơi, tài khoản hiện tại, điểm cao nhất, tổng sao, số sao đã kiếm, số sao đã dùng, số lượt chơi và tiến độ từng màn. Với mỗi màn, hệ thống lưu trạng thái hoàn thành, số sao cao nhất, điểm cao nhất và thời gian tốt nhất.

\thesistodo{nếu muốn trình bày trực quan hơn, vẽ sơ đồ cấu trúc JSON hoặc bảng mô tả các trường dữ liệu trong \texttt{records.json}.}

\begin{table}[H]
\centering
\begin{tabularx}{\textwidth}{|>{\RaggedRight\arraybackslash}p{3.8cm}|>{\RaggedRight\arraybackslash}p{3.0cm}|>{\RaggedRight\arraybackslash}X|}
\hline
\textbf{Trường dữ liệu} & \textbf{Kiểu dữ liệu} & \textbf{Ý nghĩa} \\ \hline
\texttt{current\_profile\_id} & Chuỗi & Mã tài khoản đang được chọn. \\ \hline
\texttt{profiles} & Từ điển & Danh sách tài khoản cục bộ, mỗi tài khoản có tên và bản ghi tiến độ riêng. \\ \hline
\texttt{best\_score} & Số nguyên & Điểm cao nhất của tài khoản hiện tại. \\ \hline
\texttt{total\_stars} & Số nguyên & Tổng số sao hiện có để dùng gợi ý. \\ \hline
\texttt{earned\_stars} & Số nguyên & Tổng số sao đã kiếm được từ các nhiệm vụ. \\ \hline
\texttt{spent\_stars} & Số nguyên & Tổng số sao đã sử dụng cho gợi ý. \\ \hline
\texttt{games\_played} & Số nguyên & Số màn đã hoàn thành hoặc số lượt chơi được ghi nhận. \\ \hline
\texttt{levels} & Từ điển & Tiến độ theo từng màn, khóa là mã màn chơi. \\ \hline
\texttt{completed} & Boolean & Trạng thái hoàn thành của một màn. \\ \hline
\texttt{stars} & Số nguyên & Số sao cao nhất từng đạt ở màn tương ứng. \\ \hline
\texttt{best\_time} & Số nguyên & Thời gian hoàn thành tốt nhất của màn, tính bằng giây. \\ \hline
\end{tabularx}
\caption{Các trường dữ liệu chính trong tệp lưu tiến độ}
\label{table:save_fields}
\end{table}

\section{Xây dựng ứng dụng}
\subsection{Thư viện và công cụ sử dụng}
\begin{table}[H]
\centering
\small
    \begin{tabularx}{\textwidth}{>{\RaggedRight\arraybackslash}p{3.6cm}>{\RaggedRight\arraybackslash}p{4cm}>{\RaggedRight\arraybackslash}X}
        \hline
        \textbf{Mục đích} & \textbf{Công cụ/Công nghệ} & \textbf{Ghi chú} \\ \hline
        Game engine & Godot Engine 4.x & Xây dựng giao diện và logic trò chơi \\ \hline
        Ngôn ngữ lập trình & GDScript & Viết toàn bộ logic ứng dụng \\ \hline
        Định dạng lưu trữ & JSON & Lưu tiến độ người chơi cục bộ \\ \hline
        Quản lý tài liệu & LaTeX/Overleaf & Viết báo cáo đồ án \\ \hline
        Soạn thảo mã nguồn & Visual Studio Code/Godot Editor & Phát triển và chỉnh sửa mã nguồn \\ \hline
    \end{tabularx}
    \caption{Danh sách công cụ và công nghệ sử dụng}
    \label{table:tools}
\end{table}

\thesistodo{nếu cần chính xác tuyệt đối, thay Godot 4.x bằng phiên bản cụ thể đang dùng trong Godot Editor.}

Mã nguồn phần trò chơi được tổ chức chủ yếu trong thư mục \texttt{scripts}. Tại thời điểm viết báo cáo, thư mục này gồm 13 file GDScript với khoảng 4.945 dòng mã. Các file lớn nhất là \texttt{ui\_controller.gd}, \texttt{main.gd}, \texttt{themed\_level\_provider.gd} và \texttt{manual\_clues.gd}. Điều này phản ánh đặc thù của ứng dụng: phần giao diện và dữ liệu chủ đề chiếm dung lượng đáng kể, trong khi các lớp tiện ích như chuẩn hóa tiếng Việt hoặc mô tả từ có kích thước nhỏ hơn.

\subsection{Kết quả đạt được}
Sản phẩm hiện tại là một nguyên mẫu trò chơi ô chữ tiếng Việt có thể chơi được. Ứng dụng có màn hình menu, quản lý tài khoản cục bộ, danh sách chủ đề, danh sách màn chơi, bảng ô chữ sinh tự động, gợi ý theo số, ô nhập đáp án, hệ thống điểm, hệ thống sao, nhiệm vụ hoàn thành và lưu tiến độ. Bảng chữ được hiển thị bằng các nút ô chữ, tự điều chỉnh kích thước theo vùng hiển thị để phù hợp với các kích thước màn hình khác nhau.

\begin{table}[H]
\centering
\small
    \begin{tabularx}{\textwidth}{>{\RaggedRight\arraybackslash}p{3.6cm}>{\RaggedRight\arraybackslash}p{4cm}>{\RaggedRight\arraybackslash}X}
        \hline
        \textbf{Thành phần} & \textbf{Kết quả} & \textbf{Ghi chú} \\ \hline
        Chủ đề từ vựng & Nhiều chủ đề tiếng Việt & Dữ liệu nhúng trong mã nguồn \\ \hline
        Màn chơi & Tối đa 3 màn mỗi chủ đề & Chia từ theo từng chủ đề \\ \hline
        Kích thước bảng & 18 x 18 & Cố định trong bộ sinh màn \\ \hline
        Gợi ý & Mở một chữ hoặc mở cả từ & Sử dụng sao \\ \hline
        Lưu tiến độ & Có & Theo tài khoản cục bộ \\ \hline
    \end{tabularx}
    \caption{Tóm tắt kết quả xây dựng ứng dụng}
    \label{table:results}
\end{table}

\thesistodo{thống kê số dòng mã, số file script, dung lượng mã nguồn hoặc ảnh sản phẩm đóng gói nếu giảng viên yêu cầu.}

\begin{table}[H]
\centering
\begin{tabularx}{\textwidth}{|>{\RaggedRight\arraybackslash}p{4.4cm}|>{\RaggedRight\arraybackslash}p{2.4cm}|>{\RaggedRight\arraybackslash}X|}
\hline
\textbf{Nhóm file} & \textbf{Số lượng} & \textbf{Ghi chú} \\ \hline
Script GDScript trong thư mục \texttt{scripts} & 13 & Bao gồm logic game, UI, sinh bảng, lưu dữ liệu và xử lý tiếng Việt. \\ \hline
Scene chính trong thư mục \texttt{scenes} & 4 & Gồm \classname{Main}, \classname{StartMenu}, \classname{GameBoard} và \classname{Tile}. \\ \hline
Dòng mã GDScript ước tính & 4.945 & Tính theo số dòng trong các file \texttt{.gd} hiện có. \\ \hline
Tệp cấu hình project & 1 & \texttt{project.godot} khai báo scene chạy chính và cấu hình màn hình. \\ \hline
\end{tabularx}
\caption{Thống kê sơ bộ mã nguồn ứng dụng}
\label{table:source_stats}
\end{table}

\section{Minh họa các chức năng chính}
\thesistodo{ảnh minh họa các chức năng chính. Nên có ít nhất các ảnh: màn hình chính, quản lý tài khoản, chọn chủ đề, popup chọn màn, màn chơi đang giải, màn hoàn thành.}

Các chức năng chính đã được cài đặt gồm chọn tài khoản, chọn chủ đề, chọn màn, giải từ, dùng gợi ý và xem kết quả hoàn thành. Khi người chơi chọn một gợi ý, tiêu đề gợi ý được cập nhật để hiển thị nội dung câu hỏi. Khi nhập đáp án đúng, các ô tương ứng được mở và điểm số thay đổi. Khi hoàn thành toàn bộ từ, hệ thống hiển thị số sao nhiệm vụ, điểm, thời gian và lưu kết quả.

Chức năng quản lý tài khoản cho phép nhiều người chơi sử dụng cùng một bản cài đặt mà không ghi đè tiến độ của nhau. Mỗi tài khoản có tên hiển thị và bản ghi tiến độ riêng. Khi người chơi tạo tài khoản mới, hệ thống tự tạo mã tài khoản và chuyển sang tài khoản đó. Khi xóa tài khoản, hệ thống chỉ cho phép xóa nếu còn ít nhất một tài khoản khác để tránh trạng thái không có hồ sơ.

Chức năng chọn chủ đề hiển thị tiến độ theo dạng số màn đã hoàn thành trên tổng số màn. Thông tin này giúp người chơi biết chủ đề nào đã chơi xong và chủ đề nào còn màn chưa hoàn thành. Chức năng chọn màn hiển thị số từ trong màn và số sao đã đạt, giúp người chơi quyết định chơi màn mới hoặc chơi lại màn cũ để cải thiện kết quả.

Chức năng giải ô chữ được xây dựng quanh thao tác chọn gợi ý và nhập đáp án. Người chơi không cần nhập từng ký tự vào từng ô mà nhập toàn bộ từ đoán. Cách nhập này phù hợp với tiếng Việt vì người chơi có thể dùng bộ gõ quen thuộc của hệ điều hành. Khi đáp án đúng, hệ thống tự động mở các ô chữ tương ứng trên bảng.

Chức năng gợi ý dùng sao làm tài nguyên. Người chơi có thể mở một chữ với chi phí thấp hoặc mở toàn bộ từ với chi phí cao hơn. Nếu sử dụng nhiều gợi ý, người chơi vẫn có thể hoàn thành màn nhưng có thể ảnh hưởng tới nhiệm vụ phụ. Cách thiết kế này tạo sự cân bằng giữa hỗ trợ người chơi và khuyến khích tự giải.

\section{Kiểm thử}
Quá trình kiểm thử được chia thành ba nhóm: kiểm thử chức năng, kiểm thử giao diện và trải nghiệm người dùng, và thử nghiệm chơi thực tế với nhóm nhỏ. Mục tiêu của phần này là kiểm tra mức độ ổn định và khả năng chơi được của nguyên mẫu, chưa nhằm khẳng định hiệu quả học tập hoặc khả năng cải thiện ghi nhớ từ vựng.

\subsection{Kiểm thử chức năng}
Kiểm thử chức năng tập trung vào các luồng chính của trò chơi: quản lý tài khoản, chọn chủ đề, chọn màn chơi, sinh bảng chữ, nhập đáp án, dùng gợi ý, tính điểm, tính sao và lưu tiến độ. Với sinh bảng, cần kiểm tra rằng bảng có tối thiểu số từ yêu cầu, các từ đặt trong phạm vi 18 x 18 và các từ có liên kết giao nhau. Với nhập đáp án, cần kiểm tra các trường hợp đúng dấu, có khoảng trắng, sai đáp án và bỏ trống. Với lưu tiến độ, cần kiểm tra việc tạo tài khoản mới, đổi tài khoản, hoàn thành màn và đọc lại dữ liệu sau khi khởi động lại ứng dụng.

\thesistodo{bảng test case thực tế. Mỗi dòng nên có mã test, mục tiêu, dữ liệu vào, kết quả mong đợi và kết quả thực tế. Bạn có thể bổ sung sau khi chạy game.}

\begin{table}[H]
\centering
\footnotesize
\begin{tabularx}{\textwidth}{|>{\RaggedRight\arraybackslash}p{1.05cm}|>{\RaggedRight\arraybackslash}p{2.15cm}|>{\RaggedRight\arraybackslash}X|>{\RaggedRight\arraybackslash}X|>{\RaggedRight\arraybackslash}X|>{\RaggedRight\arraybackslash}p{1.45cm}|>{\RaggedRight\arraybackslash}p{2.0cm}|}
\hline
\textbf{Mã kiểm thử} & \textbf{Chức năng} & \textbf{Dữ liệu vào} & \textbf{Kết quả mong đợi} & \textbf{Kết quả thực tế} & \textbf{Trạng thái} & \textbf{Ghi chú} \\ \hline
TC01 & Tạo tài khoản cục bộ & Nhập tên tài khoản mới và xác nhận tạo & Tài khoản được tạo, xuất hiện trong danh sách và được chọn làm tài khoản hiện tại & Ghi sau khi kiểm thử & Chưa ghi nhận & Kiểm tra cả trường hợp tên rỗng nếu có xử lý \\ \hline
TC02 & Chọn tài khoản & Chọn một tài khoản đã tồn tại & Hệ thống chuyển sang tài khoản được chọn và hiển thị tiến độ tương ứng & Ghi sau khi kiểm thử & Chưa ghi nhận & Không ghi đè tiến độ tài khoản khác \\ \hline
TC03 & Chọn chủ đề & Chọn một chủ đề trong menu & Danh sách màn của chủ đề được hiển thị đúng & Ghi sau khi kiểm thử & Chưa ghi nhận & Kiểm tra chủ đề đã/chưa hoàn thành \\ \hline
TC04 & Chọn màn chơi & Chọn một màn trong chủ đề & Hệ thống chuyển sang màn chơi, hiển thị bảng chữ và danh sách gợi ý & Ghi sau khi kiểm thử & Chưa ghi nhận & Kiểm tra màn mới và màn đã chơi \\ \hline
TC05 & Sinh bảng chữ & Tạo màn từ danh sách từ của chủ đề & Bảng được sinh trong kích thước quy định, có các từ được đặt hợp lệ & Ghi sau khi kiểm thử & Chưa ghi nhận & Kiểm tra số từ đặt được và liên kết giao nhau \\ \hline
TC06 & Chọn ô và nhập đáp án & Chọn một gợi ý, nhập đáp án vào ô nhập & Hệ thống nhận thao tác nhập và xử lý khi bấm đoán & Ghi sau khi kiểm thử & Chưa ghi nhận & Không yêu cầu nhập từng chữ vào từng ô \\ \hline
TC07 & Đáp án có dấu tiếng Việt & Nhập đúng đáp án có dấu và khoảng trắng nếu có & Đáp án được nhận đúng, từ tương ứng được mở trên bảng & Ghi sau khi kiểm thử & Chưa ghi nhận & Kiểm tra bộ gõ tiếng Việt trên môi trường chạy \\ \hline
TC08 & Đáp án không dấu & Nhập dạng không dấu của một đáp án tiếng Việt & Nếu hệ thống hỗ trợ thì chấp nhận; nếu không hỗ trợ thì báo chưa đúng nhất quán & Ghi sau khi kiểm thử & Chưa ghi nhận & Ghi rõ chính sách xử lý khi hoàn thiện tài liệu \\ \hline
TC09 & Dùng gợi ý & Chọn từ chưa hoàn thành và dùng gợi ý & Một chữ hoặc cả từ được mở theo loại gợi ý, số sao giảm đúng chi phí & Ghi sau khi kiểm thử & Chưa ghi nhận & Kiểm tra cả trường hợp không đủ sao \\ \hline
TC10 & Tính điểm & Nhập đúng/sai một số đáp án & Điểm thay đổi theo quy tắc đã cài đặt, không tăng khi đáp án sai & Ghi sau khi kiểm thử & Chưa ghi nhận & Ghi lại điểm trước và sau thao tác \\ \hline
TC11 & Tính sao & Hoàn thành màn với các điều kiện nhiệm vụ khác nhau & Số sao nhận được phản ánh kết quả nhiệm vụ của màn & Ghi sau khi kiểm thử & Chưa ghi nhận & Không kết luận cân bằng độ khó từ kiểm thử này \\ \hline
TC12 & Lưu tiến độ & Hoàn thành màn và quay về menu & Trạng thái màn, điểm/sao tốt nhất và số liệu liên quan được lưu & Ghi sau khi kiểm thử & Chưa ghi nhận & Kiểm tra theo từng tài khoản \\ \hline
TC13 & Tải lại tiến độ & Đóng và mở lại game sau khi đã lưu & Tiến độ trước đó được tải lại đúng theo tài khoản hiện tại & Ghi sau khi kiểm thử & Chưa ghi nhận & Kiểm tra tệp lưu cục bộ \\ \hline
TC14 & Đổi kích thước cửa sổ/game board & Thay đổi kích thước cửa sổ hoặc vùng hiển thị & Bảng chữ và vùng điều khiển không chồng chữ, vẫn thao tác được & Ghi sau khi kiểm thử & Chưa ghi nhận & Cần mở rộng thêm số thiết bị kiểm thử \\ \hline
TC15 & Từ không đặt được vào bảng & Dùng danh sách có từ dài hoặc khó giao nhau & Hệ thống bỏ qua hoặc xử lý từ không đặt được mà không treo game & Ghi sau khi kiểm thử & Chưa ghi nhận & Ghi số từ đặt được nếu có thống kê \\ \hline
\end{tabularx}
\caption{Khung kiểm thử chức năng}
\label{table:test_cases}
\end{table}

\subsection{Kiểm thử giao diện và trải nghiệm người dùng}
Kiểm thử giao diện tập trung vào khả năng đọc, khả năng thao tác và phản hồi của các màn hình chính. Các điểm cần quan sát gồm: menu chọn tài khoản có dễ hiểu hay không, luồng chọn chủ đề và chọn màn có rõ ràng hay không, bảng chữ có bị chồng chữ khi thay đổi kích thước cửa sổ hay không, trạng thái đúng/sai/gợi ý có phản hồi kịp thời hay không, và hộp thoại hoàn thành có cung cấp đủ thông tin cần thiết hay không.

\begin{table}[H]
\centering
\footnotesize
\begin{tabularx}{\textwidth}{|>{\RaggedRight\arraybackslash}p{2.8cm}|>{\RaggedRight\arraybackslash}X|>{\RaggedRight\arraybackslash}X|>{\RaggedRight\arraybackslash}p{2.2cm}|}
\hline
\textbf{Nhóm kiểm tra} & \textbf{Tiêu chí quan sát} & \textbf{Cách kiểm tra} & \textbf{Ghi nhận} \\ \hline
Menu và tài khoản & Người chơi nhận biết được tài khoản hiện tại và thao tác tạo/chọn tài khoản & Quan sát người chơi đi từ màn đầu tới bước chọn tài khoản & Ghi sau khi kiểm thử \\ \hline
Chủ đề và màn chơi & Người chơi hiểu được tiến độ chủ đề và biết cách vào một màn & Yêu cầu người chơi chọn một chủ đề và một màn cụ thể & Ghi sau khi kiểm thử \\ \hline
Bảng chữ & Ô chữ, gợi ý và vùng nhập không chồng lên nhau ở các kích thước phổ biến & Thử ở ít nhất hai kích thước cửa sổ khác nhau & Ghi sau khi kiểm thử \\ \hline
Phản hồi thao tác & Trạng thái đúng, sai, thiếu sao hoặc hoàn thành được hiển thị rõ & Nhập đúng, nhập sai, dùng gợi ý và hoàn thành màn & Ghi sau khi kiểm thử \\ \hline
Mức độ dễ dùng & Người chơi có thể hoàn thành luồng chơi cơ bản mà không cần hướng dẫn quá nhiều & Ghi số lần người chơi hỏi lại hoặc thao tác nhầm & Ghi sau khi kiểm thử \\ \hline
\end{tabularx}
\caption{Các tiêu chí kiểm thử giao diện và trải nghiệm người dùng}
\label{table:ui_ux_test_criteria}
\end{table}

\subsection{Thử nghiệm chơi thực tế với nhóm nhỏ}
Thử nghiệm chơi thực tế được định hướng thực hiện sơ bộ với một nhóm nhỏ gồm bạn bè hoặc người quen. Đây là thử nghiệm sơ bộ với nhóm nhỏ, không chọn mẫu ngẫu nhiên, không đại diện cho toàn bộ người dùng. Kết quả, khi được đo thực tế, chỉ nên dùng để đánh giá sơ bộ khả năng chơi được, độ ổn định của luồng chức năng và các vấn đề giao diện dễ nhận thấy.

Nếu tỉ lệ qua màn ghi nhận ở mức cao, cần diễn giải thận trọng vì người chơi quen người phát triển, số lượng màn thử chưa nhiều, người chơi có thể được hỗ trợ trong lúc thử, và chưa có nhóm người dùng đại diện. Vì vậy, không sử dụng các số liệu này để kết luận trò chơi có hiệu quả học tập.

\begin{table}[H]
\centering
\footnotesize
\begin{tabularx}{\textwidth}{|>{\RaggedRight\arraybackslash}p{3.0cm}|>{\RaggedRight\arraybackslash}X|>{\RaggedRight\arraybackslash}X|}
\hline
\textbf{Nội dung ghi nhận} & \textbf{Cách ghi nhận} & \textbf{Lưu ý khi diễn giải} \\ \hline
Tỉ lệ qua màn & Ghi số màn người chơi thử và số màn hoàn thành. & Chỉ phản ánh nhóm nhỏ, không đại diện cho toàn bộ người dùng. \\ \hline
Thời gian hoàn thành & Ghi thời gian trung bình theo từng màn hoặc từng lượt chơi. & Không dùng để kết luận độ khó tổng quát nếu số màn thử còn ít. \\ \hline
Số lần dùng gợi ý & Ghi số lượt mở chữ hoặc mở cả từ trong mỗi màn. & Cần xem cùng độ khó của từ và mức quen thuộc của người chơi. \\ \hline
Nhận xét giao diện & Ghi góp ý về kích thước chữ, vùng nhập, nút bấm và phản hồi đúng/sai. & Người quen có thể đưa phản hồi tích cực hơn nhóm người dùng độc lập. \\ \hline
\end{tabularx}
\caption{Các nội dung cần ghi nhận khi thử nghiệm chơi thực tế với nhóm nhỏ}
\label{table:small_playtest_results}
\end{table}

\thesistodo{bổ sung bảng kết quả thử nghiệm thực tế sau khi có số liệu đo thật; không dùng số liệu minh họa như dữ liệu thực nghiệm.}

\subsection{Hạn chế của kiểm thử}
Phần kiểm thử hiện tại còn một số hạn chế. Số lượng người tham gia còn ít và người tham gia chủ yếu là bạn bè hoặc người quen nên có thể có thiên lệch tích cực. Quá trình thử nghiệm chưa được thực hiện trên nhiều thiết bị và nhiều kích thước màn hình khác nhau. Báo cáo cũng chưa có đánh giá dài hạn về khả năng ghi nhớ từ vựng sau khi chơi. Do đó, kết quả kiểm thử chỉ phản ánh mức độ ổn định và khả năng chơi được của nguyên mẫu tại thời điểm thử nghiệm, chưa đủ cơ sở để khẳng định hiệu quả học tập.

\section{Triển khai}
Ứng dụng được triển khai thử nghiệm ở dạng project Godot chạy cục bộ trên máy tính. Người dùng mở project bằng Godot hoặc chạy bản export nếu đã đóng gói. Dữ liệu tiến độ được lưu trong vùng \texttt{user://} của Godot nên không cần máy chủ hoặc kết nối mạng.

\thesistodo{môi trường triển khai cụ thể, ví dụ Windows 10/11, độ phân giải màn hình, phiên bản Godot và nếu có thì đường dẫn/bản build đã export.}

Do ứng dụng hiện là trò chơi một người chơi cục bộ, quá trình triển khai đơn giản hơn so với các hệ thống có máy chủ. Khi xuất bản cho Windows, cần cấu hình export preset trong Godot, chọn biểu tượng nếu có, sau đó tạo tệp thực thi. Khi chạy lần đầu, ứng dụng sẽ tự tạo tệp lưu nếu chưa có. Nếu muốn phát hành bản Web, cần kiểm tra lại kích thước giao diện và khả năng nhập tiếng Việt trong trình duyệt, vì bộ gõ tiếng Việt có thể có hành vi khác giữa môi trường desktop và web.

\end{document}
